
\documentclass{article}
\usepackage{colortbl}
\usepackage{makecell}
\usepackage{multirow}
\usepackage{supertabular}

\begin{document}

\newcounter{utterance}

\centering \large Interaction Transcript for game `wordle', experiment `medium\_frequency\_words\_no\_clue\_no\_critic', episode 11 with reflexion{-}Qwen3.5{-}27B.
\vspace{24pt}

{ \footnotesize  \setcounter{utterance}{1}
\setlength{\tabcolsep}{0pt}
\begin{supertabular}{c@{$\;$}|p{.15\linewidth}@{}p{.15\linewidth}p{.15\linewidth}p{.15\linewidth}p{.15\linewidth}p{.15\linewidth}}
    \# & \multicolumn{2}{c}{Player} && \multicolumn{2}{c}{Game Master} \\
    \hline

    \theutterance \stepcounter{utterance}  
    & & & \multicolumn{4}{p{0.6\linewidth}}{
        \cellcolor[rgb]{0.9,0.9,0.9}{
            \makecell[{{p{\linewidth}}}]{
                \texttt{\tiny{[P1$\langle$GM]}}
                \texttt{You have attempted this same task before and did not succeed. These are the plans you wrote after those attempts.} \\
\texttt{Trial \#1: New plan: The previous attempt failed because I included my internal reasoning process (e.g., "wait, 'a' was yellow...", "Let's try...") directly in the "explanation" field, causing the response to become disorganized and likely violating the "concise explanation" rule or the strict format. Additionally, I did not provide a final "guess:" line before the text cut off. Next time, I will strictly separate my internal logic from the output. I will perform the deduction mentally to identify the correct word (likely "taint" or "tania" based on the constraints: 'a' not in 1,2,3; 't' not in 4; 'i' not in 2; 'n' not in 3), and then output *only* a clean, concise explanation of the final choice followed immediately by the "guess:" tag and the five{-}letter word. I will ensure the explanation summarizes the logic without showing the step{-}by{-}step trial{-}and{-}error or self{-}correction.} \\
\\ 
\\ 
\texttt{You are a language wizard who likes to guess words by using the given rules.} \\
\\ 
\texttt{Welcome to Wordle! You have six attempts to guess the target word, a valid English word of five lowercase letters (a{-}z). Please use the tags "explanation:" and "guess:" to provide a concise explanation for each guess.} \\
\\ 
\texttt{For instance, if your guess is "apple", your response should be} \\
\texttt{explanation: this is a common five{-}letter English word, and I am starting my guess with this word.} \\
\texttt{guess: apple} \\
\\ 
\texttt{After each guess, your answer will be validated, and you will receive feedback indicating which letters are correct (green), which letters are correct but in the wrong position (yellow), and which letters are incorrect (red). This feedback can be useful in determining which letters to include or exclude in your next guess.} \\
\\ 
\texttt{For example, the feedback for "apple" might be:} \\
\texttt{guess\_feedback: a<yellow> p<yellow> p<green> l<yellow> e<red>} \\
\\ 
\texttt{The explanation should contain details about how the guess\_feedback is used to arrive at a new guess.} \\
\\ 
\texttt{Do not generate any other text.} \\
\\ 
\texttt{Let's begin with your first guess.} \\
            }
        }
    }
     \\ \\

    \theutterance \stepcounter{utterance}  
    & \multicolumn{4}{p{0.6\linewidth}}{
        \cellcolor[rgb]{0.561,0.737,0.561}{
            \makecell[{{p{\linewidth}}}]{
                \texttt{\tiny{[P1$\rangle$GM]}}
                \texttt{explanation: I am starting with a common five{-}letter word containing frequent vowels and consonants to gather initial feedback on letter presence and positions.} \\
\texttt{guess: slate} \\
            }
        }
    }
    & & \\ \\

    \theutterance \stepcounter{utterance}  
    & & & \multicolumn{2}{p{0.3\linewidth}}{
        \cellcolor[rgb]{0.9,0.9,0.9}{
            \makecell[{{p{\linewidth}}}]{
                \texttt{\tiny{[GM$|$GM]}}
                \texttt{attempts = 1} \\
\texttt{target\_word = faint} \\
\texttt{guess = slate} \\
\texttt{guess\_feedback = s<red> l<red> a<yellow> t<yellow> e<red>} \\
            }
        }
    }
    & & \\ \\

    \theutterance \stepcounter{utterance}  
    & & & \multicolumn{4}{p{0.6\linewidth}}{
        \cellcolor[rgb]{0.9,0.9,0.9}{
            \makecell[{{p{\linewidth}}}]{
                \texttt{\tiny{[P1$\langle$GM]}}
                \texttt{guess\_feedback: s<red> l<red> a<yellow> t<yellow> e<red>} \\
\\ 
\texttt{Provide your response only in this format.} \\
\texttt{explanation: details} \\
\texttt{guess: word} \\
            }
        }
    }
     \\ \\

    \theutterance \stepcounter{utterance}  
    & \multicolumn{4}{p{0.6\linewidth}}{
        \cellcolor[rgb]{0.561,0.737,0.561}{
            \makecell[{{p{\linewidth}}}]{
                \texttt{\tiny{[P1$\rangle$GM]}}
                \texttt{explanation: The feedback indicates 'a' and 't' are in the word but not in positions 3 and 4 respectively, while 's', 'l', and 'e' are absent. I will try a word that places 'a' and 't' in new positions and uses common remaining letters like 'r', 'i', and 'n'.} \\
\texttt{guess: train} \\
            }
        }
    }
    & & \\ \\

    \theutterance \stepcounter{utterance}  
    & & & \multicolumn{2}{p{0.3\linewidth}}{
        \cellcolor[rgb]{0.9,0.9,0.9}{
            \makecell[{{p{\linewidth}}}]{
                \texttt{\tiny{[GM$|$GM]}}
                \texttt{attempts = 2} \\
\texttt{target\_word = faint} \\
\texttt{guess = train} \\
\texttt{guess\_feedback = t<yellow> r<red> a<yellow> i<yellow> n<yellow>} \\
            }
        }
    }
    & & \\ \\

    \theutterance \stepcounter{utterance}  
    & & & \multicolumn{4}{p{0.6\linewidth}}{
        \cellcolor[rgb]{0.9,0.9,0.9}{
            \makecell[{{p{\linewidth}}}]{
                \texttt{\tiny{[P1$\langle$GM]}}
                \texttt{guess\_feedback: t<yellow> r<red> a<yellow> i<yellow> n<yellow>} \\
\\ 
\texttt{Provide your response only in this format.} \\
\texttt{explanation: details} \\
\texttt{guess: word} \\
            }
        }
    }
     \\ \\

    \theutterance \stepcounter{utterance}  
    & \multicolumn{4}{p{0.6\linewidth}}{
        \cellcolor[rgb]{0.561,0.737,0.561}{
            \makecell[{{p{\linewidth}}}]{
                \texttt{\tiny{[P1$\rangle$GM]}}
                \texttt{explanation: The feedback shows 't', 'a', 'i', and 'n' are present but misplaced, while 'r' is absent. 't' is not in position 1, 'a' not in 3, 'i' not in 4, and 'n' not in 5. The word "taint" places 't' in position 1 (incorrect), so I need to move 't'. "tania" is not a common word. "taint" has 't' at 1, 'a' at 2, 'i' at 3, 'n' at 4, 't' at 5. Wait, 't' was yellow in pos 1, so it can't be pos 1. 'a' was yellow in pos 3, so it can't be pos 3. 'i' was yellow in pos 4, so it can't be pos 4. 'n' was yellow in pos 5, so it can't be pos 5. The letters are t, a, i, n, and one unknown. Let's try "taint" again? No, 't' at 1 is wrong. How about "titan"? 't' at 1 is wrong. "taint" is t{-}a{-}i{-}n{-}t. 't' at 1 (bad), 'a' at 2 (ok), 'i' at 3 (ok), 'n' at 4 (ok), 't' at 5 (ok). But 't' was yellow at 1, so it must be elsewhere. If the word is "taint", the first 't' is wrong position, second 't' is right? No, 't' is yellow in pos 1, meaning 't' is in the word but not at 1. If the word is "taint", 't' is at 1 and 5. If 't' is at 5, that's fine. But 'a' was yellow at 3, so 'a' cannot be at 3. In "taint", 'a' is at 2. 'i' was yellow at 4, so 'i' cannot be at 4. In "taint", 'i' is at 3. 'n' was yellow at 5, so 'n' cannot be at 5. In "taint", 'n' is at 4. This fits. But wait, 't' is at 1 in "taint", which is the position that was yellow. So 't' cannot be at 1. So "taint" is invalid because of the first 't'. What about "tania"? Not a word. "taint" is the only common word with these letters? Maybe "titan"? 't} \\
            }
        }
    }
    & & \\ \\

    \theutterance \stepcounter{utterance}  
    & & & \multicolumn{2}{p{0.3\linewidth}}{
        \cellcolor[rgb]{0.9,0.9,0.9}{
            \makecell[{{p{\linewidth}}}]{
                \texttt{\tiny{[GM$|$GM]}}
                \texttt{Error: The response should contain the 'guess:' keyword exactly once.} \\
            }
        }
    }
    & & \\ \\

    \theutterance \stepcounter{utterance}  
    & & & \multicolumn{2}{p{0.3\linewidth}}{
        \cellcolor[rgb]{0.9,0.9,0.9}{
            \makecell[{{p{\linewidth}}}]{
                \texttt{\tiny{[GM$|$GM]}}
                \texttt{game\_result = ABORT} \\
            }
        }
    }
    & & \\ \\

    \theutterance \stepcounter{utterance}  
    & & \multicolumn{4}{p{0.6\linewidth}}{
        \cellcolor[rgb]{0.95,0.95,0.95}{
            \makecell[{{p{\linewidth}}}]{
                \texttt{\tiny{[P1$\leftrightarrow$P]}}
                \texttt{Post{-}episode Reflexion (not sent to GM):} \\
\texttt{New plan: The previous attempt failed because I included my internal trial{-}and{-}error reasoning (e.g., "Wait...", "Let's try...", "No...") within the "explanation" field, violating the requirement for a concise explanation, and I failed to output the final "guess:" line before the response was cut off. Next time, I will perform all logical deductions and word validation internally before generating the response. I will ensure the "explanation" field contains only a clear, final summary of why the chosen word fits the constraints (e.g., "Based on feedback, 't', 'a', 'i', 'n' are present but misplaced; 'taint' places them in valid positions while avoiding excluded spots"). I will strictly adhere to the format by providing the explanation followed immediately by a valid five{-}letter word on the "guess:" line, ensuring no internal monologue or incomplete sentences appear in the output.} \\
            }
        }
    }
    & \\ \\

\end{supertabular}
}

\end{document}
